
\section{Mechanization}
\label{sec:mechanization}

Everything this paper states is mechanized in Rocq: the calculus of
\cref{sec:calculus}, the logical relation of \cref{sec:logicalRelation}, the
extensions of \cref{sec:extensions}, and all examples. The development is built
on top of the \texttt{std++} library and does not admit any axiom.

The conception of the semantics, type system, logical relation, examples and sketches of the soundness arguments were all done without the use of AI.
The mechanization was produced with the aid of generative AI
(Claude Code, Fable 5 orchestrating Opus 5),
in several refinements starting from the manually written definitions and proof sketches.
We thoroughly audited the definitions for accuracy against their hand-written versions.
\ed{\pt! \mw! please check you agree with this characterization}

The artifact comprises roughly $17{,}300$ lines (non-blank, non-comment) of
Rocq across 30 files, containing about $300$ definitions and $850$ lemmas and
theorems.
% \Cref{tab:mechanization} shows the distribution of LOC over
% components of the formalization.
\ed{To be honest I don't think the table adds useful information. I think the interesting aspect is how much of it is examples.}
A significant portion of the code ($\sim40\%$) is spent fully working out examples.
The source of complexity is the semantic typing proofs, which need to reason
about the LTS directly. To tame the combinatorial explosion, the proofs use a simple invariant-based method. The semantic proofs are still much harder than they need to be, as is to be expected, which would call for using a program logic layer to perform them in a more scalable way. We leave this exploration for future work.

% \begin{table}[t]
%   \caption{Lines of Rocq by component.}
%   \label{tab:mechanization}
%   \begin{tabular}{lr}
%     \toprule
%     Component                                                                                 & LOC          \\
%     \midrule
%     Syntax, types, syntactic typing, substitution boilerplate (\cref{fig:syntax,fig:typeSys}) & 1,829        \\
%     Labelled semantics and harmony (\cref{fig:term-semantics,fig:proc-semantics})             & 618          \\
%     The logical relation (\cref{fig:logrel-rules-term,fig:logrel-rules-proc})                 & 642          \\
%     Safety and adequacy                                                                       & 442          \\
%     Compositionality for terms                                                                & 1,117        \\
%     Compositionality for processes                                                            & 3,003        \\
%     Other                                                                                     & 2,643        \\
%     Examples                                                                                  & 7,025        \\
%     \midrule
%     Total                                                                                     & $\sim$17,300 \\
%     \bottomrule
%   \end{tabular}
% \end{table}

% The main proof work is done in \ref{rule:P-Connect} (\cref{lem:comp:connect}),
% which alone accounts for $2{,}041$ lines of the $3{,}003$ for the
% compositionality theorem of processes. Connecting two processes over a channel
% is the linear-logic cut, and it is the only place where membership in the
% logical relation for the combined process is hard to maintain.

% We also spend a lot of code to formalize examples, nearly $41\%$ of the
% development. This is not incidental, since a process is admitted by the
% relation only if every transition it can take discharges its obligation, so an
% example is a proof about a whole reachable state space. The load balancer of
% \cref{ex:load-balancer} takes $2{,}270$ lines and the leader election of
% \cref{sec:extension:-sums} $1{,}981$ lines to conclude membership in the
% logical relation.

% \mw{TODO: explain invariant method @emanuele}

%%% Local Variables:
%%% mode: latex
%%% TeX-master: "main"
%%% End:
